\section{Related Work}
\label{sec:rw}

This section reviews the state of the art in cloud autoscaling, organized into five thematic areas: reactive autoscaling mechanisms, proactive workload forecasting, reinforcement learning approaches, stability and oscillation mitigation, and the specific gap this project addresses.

\subsection{Reactive Autoscaling in Kubernetes}

Kubernetes' built-in Horizontal Pod Autoscaler (HPA) and event-driven autoscalers like KEDA are both reactive: they scale only after a metric threshold is already crossed, which under bursty workloads leads to either SLA-violating under-provisioning or costly over-provisioning. The HPA control loop has been formally analyzed for stability properties by recent work \cite{HpaStability25}, confirming that threshold-based reactive scaling converges to steady state but suffers from inherent reaction latency during traffic spikes.

Several comprehensive surveys document the limitations of reactive approaches. A 2024 systematic review \cite{AutoScalingSurvey24} categorizes autoscaling techniques by machine learning paradigm (supervised, unsupervised, reinforcement learning, semi-supervised) and notes that reactive mechanisms ``trigger scaling actions only after resource utilization exceeds predefined thresholds,'' causing cold-start delays, over-provisioning, under-provisioning, and SLO violations in dynamic multi-tenant environments. The review identifies ML-based proactive forecasting as a key research direction to overcome these limitations.

\subsection{Proactive Workload Forecasting}

To address reactive limitations, researchers have developed predictive autoscalers using time-series forecasting. LSTM-based models are widely adopted for cloud workload prediction due to their ability to capture temporal dependencies. Rossi et al.~\cite{RossiUncertaintyAware24} present uncertainty-aware LSTM forecasting with transfer learning, demonstrating that probabilistic predictions allow cloud managers to make informed decisions by adjusting safety margins based on prediction confidence. Their work on Google and Alibaba traces highlights the importance of epistemic and aleatoric uncertainty quantification.

More recent deep learning architectures have pushed prediction accuracy further. Arbat et al.~\cite{ArbatWassersteinTransformer22} propose a Wasserstein Adversarial Transformer (WGAN-gp Transformer) for VM auto-scaling that achieves 5× faster inference and 5.1\% higher accuracy than state-of-the-art LSTM models. When integrated into Google Cloud Platform autoscaling, the Transformer-based mechanism significantly reduces both over-provisioning and under-provisioning rates. Similarly, transformer-based workload prediction integrated with adaptive autoscaling mechanisms has been shown to reduce SLA violations by 28\% and compute costs by 16\% under bursty conditions \cite{TransformerAutoScaling26}.

Deep CNN-LSTM hybrid architectures have also proven effective. A DCNN-LSTM model \cite{DeepCNNLSTM24} for CPU consumption prediction and SLA violation detection outperforms ARIMA, standalone LSTM, and CNN baselines, achieving 22.4\% better performance than ARIMA on the composite Energy-SLA-Violation metric. Multivariate time-series forecasting using bidirectional LSTM (BiLSTM) has been applied to predict multiple resource parameters simultaneously (CPU, memory, disk I/O, network throughput) \cite{BiLSTMWorkload23}, recognizing that one-dimensional predictions fail to capture the relationship between different resource requirements. A multiplicative LSTM variant (mLSTM) \cite{mLSTMWorkload24} addresses long-term dependencies in workload prediction, outperforming standard LSTM variants in prediction accuracy for reducing SLA violations.

Load-aware predictive frameworks that forecast both future workloads and the cloud resources necessary to cope with demand have been extensively tested in simulation \cite{LoadAwarePredictive26}, showcasing effectiveness in scaling resources according to predicted arrival patterns while maintaining desired utilization levels. PAHPA (Prediction-Assisted HPA) \cite{PAHPA26} integrates a self-updating LightGBM prediction model with queueing theory analysis and real-time correction mechanisms, achieving 9.5\% lower violation rates (16.3\% vs. 25.8\% for HPA) and 63\% reduction in 99th percentile latency. A key insight from PAHPA is that ``predictions should not be the sole basis for decisions and must be complemented by real-time monitoring data.''

\subsection{Reinforcement Learning for Cloud Autoscaling}

Reinforcement learning, particularly Deep Q-Networks (DQN), has emerged as a promising approach for adaptive autoscaling policies that learn optimal resource allocation strategies through interaction with the environment. A 2024 study \cite{OptimizingRL24} evaluating DQN, Proximal Policy Optimization (PPO), and Actor-Critic methods for cloud resource allocation found that RL approaches achieve average resource usage efficiencies of 85.7\% (DQN), 88.1\% (PPO), and 87.4\% (Actor-Critic), significantly outperforming traditional rule-based systems (72.3\%). The RL techniques reduced costs by up to 29.4\%.

DQN-based approaches have been applied to multi-dimensional autoscaling problems. A dual scheduling framework using two layers of DQN algorithms \cite{DualSchedulingDQN25} addresses both QoS-aware task allocation and workload-aware VM auto-scaling for SaaS providers, achieving higher task success rates, improved resource utilization, and significantly reduced VM rental costs compared to baseline solutions. Multi-dimensional autoscaling in microservices using DQN \cite{MultiDimDQN25}, capable of adjusting both pod count and CPU allocation simultaneously, demonstrates improved responsiveness, reduced response-time spikes, and more efficient resource use compared to single-dimension RL autoscalers and standard HPA.

In multi-cloud environments, DRL-based dynamic resource allocation frameworks using PPO and DQN \cite{MultiCloudDRL25} enable intelligent workload distribution that maximizes system efficiency, lowers latency and resource waste, and outperforms heuristic-based and rule-based approaches in terms of response time, scalability, and cost-effectiveness. For multi-tenant clouds, a DQN framework combined with predictive modeling \cite{DeepRLMultiTenant25} reduces over-provisioning and SLA violations while saving costs, with experiments on Microsoft Azure and MIT Supercloud datasets confirming effectiveness.

DRL has also been applied to specialized domains. Deep recurrent-reinforcement learning for serverless function autoscaling \cite{DeepRecurrentRL24} and RL-based vertical memory scaling for HPC workloads in Kubernetes (VERA) \cite{VERA26} demonstrate that observation-driven RL recommenders can outperform retrospective heuristics, reclaiming significantly more memory headroom while avoiding out-of-memory events.

\subsection{Binding baseline: Kubernetes HPA}

The binding CA2 baseline is traditional Kubernetes Horizontal Pod Autoscaler (CPU utilization, \texttt{autoscaling/v2}), not NimbusGuard. HPA is reactive: it scales after a metric already exceeds a target. Formal comparison is PAKS (forecast + Kubernetes API scale) versus that HPA control loop.

\subsection{Related: Proactive Autoscaling with Deep Q-Networks}

Wanigasooriya \& Ekanayake \cite{Wanigasooriya26}, ``NimbusGuard: A Novel Framework for Proactive Kubernetes Autoscaling Using Deep Q-Networks'' (IEEE ICOIN 2026), is \textbf{related literature}. NimbusGuard combines a Deep Q-Network scaling agent with an LSTM workload forecaster, evaluated against HPA and KEDA. Their results name an agility--stability trade-off. A superseded proxy in this repository reproduced that niche with an MLP simulator; it is not the binding CA2 contract.

\subsection{Stability, Oscillation Mitigation, and Hysteresis}

While proactive forecasting improves SLA compliance, it can introduce instability and oscillation (rapid, unnecessary scaling changes) that waste resources. Oscillation mitigation has received insufficient attention in the literature \cite{OscillationMitigation24}, but several techniques have emerged.

Smoothing and hysteresis mechanisms are fundamental to autoscaler stability. A theoretical analysis of systems with multiple operating levels \cite{MultipleOperatingLevels18} formalizes the use of exponential smoothing to avoid overreacting to spurious workload changes and hysteresis (different thresholds for increasing vs. decreasing levels) to avoid rapid oscillations. This work models cloud-like systems with activation delays and demonstrates that ``inertia'' through averaging and ``hysteresis'' through threshold separation are essential for stability.

Kubernetes HPA itself incorporates a downscale stabilization window (default 5 minutes) that records all scale recommendations and selects the highest within the window before scaling down, smoothing the impact of rapidly fluctuating metrics. The official documentation warns that setting the cooldown delay too long makes HPA unresponsive to workload changes, while setting it too short allows thrashing \cite{K8sHPADoc}. Automatic data featurization with grid-based oscillation mitigation \cite{GridOscillation24} dynamically matches scaling actions to reference actions using historical data, offering an alternative to fixed cooldown timers that require optimized delay values.

Recent work on attention-enhanced LSTM for Kubernetes autoscaling \cite{AttentionLSTM26} addresses ``temporal blindness'' in DRL agents (inability to capture long-term dependencies) by employing a deep temporal attention mechanism that selectively weights historical states, filtering high-frequency noise while retaining critical demand-shift precursors. This approach reduces 90th percentile latency by 29\% and replica churn by 39\% compared to single-layer LSTM baselines, confirming that ``mitigating temporal blindness through deep attentive memory is a prerequisite for reliable, low-jitter autoscaling.''

A comprehensive Kubernetes autoscaling review \cite{K8sAutoscalingReview26} summarizes that supervised time-series forecasting is the most mature direction for proactive HPA, graph-based learning is increasingly important for dependency-aware microservice autoscaling, and reinforcement learning is promising for adaptive closed-loop control, but notes that stability mechanisms remain underexplored.

\subsection{The Gap This Project Targets}

Formal CA2 targets the gap between reactive HPA and a predictive, Kubernetes-API-driven scaler (PAKS) under trace-driven LSTM forecasts. Literature on LSTM/HPA hybrids and on NimbusGuard-style RL agents informs design but does not replace the formal datasets (GCT/Alibaba), AWS testbed, or metric suite. The proxy MLP experiment addresses a narrower stability trade-off on synthetic load and is labelled as such.

